\chapter{CLI и типовые сценарии}

\section{Главная команда}

Для большинства пользовательских запусков следует начинать с \cli{process-archive-path}. Без \cli{--execute} она выполняет planning/preflight, с \cli{--execute} --- записывает результаты и запускает pipeline.

\begin{lstlisting}
python -m grandgousier.cli process-archive-path `
  --path "D:\archive\series" `
  --execute `
  --manifest-out reports\series.json
\end{lstlisting}

\section{Явный существующий каталог}

Для воспроизводимых повторных запусков можно передавать заранее известный identity catalog:

\begin{lstlisting}
python -m grandgousier.cli process-archive-path `
  --path "D:\archive\series" `
  --execute `
  --catalog-name "existing_catalog_name"
\end{lstlisting}

Пустой explicit catalog считается precondition error, если пользователь явно не разрешил отладочный режим пустого каталога.

\section{Selective rerun существующей series}

Начиная с 0.5.45.1 \cli{process-archive-path} может повторно запускать только явно выбранную
существующую series, не перестраивая field-series grouping и не перезапуская соседние series из
того же каталога:

\begin{lstlisting}
python -m grandgousier.cli process-archive-path `
  --path "D:\archive\night" `
  --execute `
  --series-id 78 `
  --catalog-name "existing_catalog_name" `
  --ensemble-min-ref-stars-per-frame 20
\end{lstlisting}

Параметр \cli{--series-id} можно повторять. В selective mode программа проверяет, что каждая
запрошенная series действительно содержит кадры, найденные под \cli{--path}; иначе запуск
останавливается с \path{FAILED_PRECONDITION_SERIES_SELECTION}. Новые series в этом режиме
не создаются. Default conventional ensemble threshold остаётся 30; меньшее значение следует
задавать только по диагностически подтверждённой причине.

\section{Color calibration}

\begin{lstlisting}
--enable-color-calibration
\end{lstlisting}

При штатной политике reference photometry программа может автоматически получить недостающее покрытие внешнего каталога. Для полностью локального запуска используйте соответствующий запрет automatic reference fetch.

\section{Candidate-excluded ensemble}

Режим до сходимости:

\begin{lstlisting}
--exclude-candidates-from-ensemble `
--candidate-ensemble-until-converged `
--candidate-ensemble-convergence-max-iterations 10 `
--candidate-ensemble-growth-policy initial_only
\end{lstlisting}

Политика \cli{initial\_only} фиксирует исходный exclusion set и не позволяет новым diagnostics каскадно расширять reference exclusions.


\section{Удалённый MASTER archive: discovery, download и reduction}

Stage~1A отделяет поиск удалённых наблюдений от реальной передачи файлов. Discovery поддерживает \cli{auto|postgres|http}; HTTPS mode предназначен в том числе для рабочих станций, которым MASTER PostgreSQL закрыт через \cli{pg\_hba.conf}. Начиная с 0.5.41.7 HTTPS discovery использует штатный параметр \cli{coord2000} веб-интерфейса и локально уточняет область до запрошенного \cli{--radius-deg}; допустимый радиус для HTTPS сейчас не больше $0.95^\circ$. Сначала создайте inventory:

\begin{lstlisting}
python -m grandgousier.cli master-archive-search `
  --search-transport http `
  --http-user "googlov" `
  --ra 327.60625 `
  --dec 72.1686111 `
  --radius-deg 0.8 `
  --out "C:\master\gpx\inventory.csv"
\end{lstlisting}

Затем проверьте download plan без \cli{--execute}:

\begin{lstlisting}
python -m grandgousier.cli master-archive-fetch `
  --manifest "C:\master\gpx\inventory.csv" `
  --out-root "C:\master\gpx"
\end{lstlisting}

Для реального transfer добавьте \cli{--execute}. Dark/flat дедуплицируются по remote path, существующие файлы по умолчанию не перезаписываются. Для non-interactive запуска задайте \cli{GRANDGOUSIER\_MASTER\_HTTP\_USER} и \cli{GRANDGOUSIER\_MASTER\_HTTP\_PASSWORD}; при отсутствии credentials команда завершается fail-fast и не открывает prompt.

HTTP downloader проверяет authoritative размер из \cli{Content-Length} или \cli{Content-Range}, не принимает resume с несовпадающим началом Range и атомарно публикует файл только после проверки полного размера. При временной ошибке сохраняется \cli{.part}; следующая попытка возобновляет передачу с его фактического размера. Прогресс передаётся ограниченно по времени (по умолчанию каждые 15 секунд), а отсутствие реального прироста байтов в течение 300 секунд прерывает попытку и оставляет \cli{.part} для retry. Размер stall timeout и интервал можно изменить в секции \cli{archives.*.assets}: \cli{stall\_timeout\_seconds} и \cli{progress\_interval\_seconds}.


После загрузки сначала проверьте Stage~1B plan:

\begin{lstlisting}
python -m grandgousier.cli master-archive-reduce `
  --manifest "C:\master\gpx\inventory.csv" `
  --cache-root "C:\master\gpx"
\end{lstlisting}

Для записи processed FITS добавьте \cli{--execute}. По умолчанию результаты появляются в \cli{<cache-root>/processed/YYYY/MM/DD/<raw\_id>\_processed.fits}; одновременно пишется \cli{master\_reduction\_report.csv}. Редуктор не масштабирует SDARK по экспозиции и не перенормирует SFLAT: несовместимые calibration frames считаются precondition error.

\section{Диагностические параметры}

Performance/diagnostic flags предназначены для проверки алгоритма и обычно не нужны в массовом production-run. Их следует включать только при исследовании конкретной стадии и затем возвращаться к обычной команде.
